\documentclass[11pt]{article}
\usepackage[margin=1in]{geometry}
\usepackage{amsmath,amssymb,amsthm,mathtools}
\usepackage{graphicx}
\usepackage{microtype}
\usepackage[colorlinks=true,linkcolor=blue,citecolor=blue,urlcolor=blue]{hyperref}

\newtheorem{theorem}{Theorem}
\newtheorem{remark}{Remark}
\newcommand{\A}{A^{**}}
\newcommand{\Z}{\mathcal Z}

\title{The General Converses for Two-Sided Orthogonality Maps}
\author{Run \texttt{fa\_banach\_001}, lane 0
  (\texttt{agent\_lane\_00})}
\date{17 August 2026}

\begin{document}
\maketitle

\begin{abstract}
Fadaee and Ghahramani ask whether the converses of their Results 3.2 and
3.3 hold for general $C^*$-algebras.  We give direct affirmative proofs.
If $\delta=d+\eta(\cdot)$ with $d$ a derivation and
$\eta\in\Z(A^{**})$, then the two zero-product identities are just the
Leibniz rule plus central terms proportional to $ab$ or $ba$.  If
$\delta=d_1+\eta(\cdot)=d_2+(\cdot)\eta$ with $d_1,d_2$ star derivations,
use the left representation for the identity at $ab^*=0$ and the right
representation for the identity at $b^*a=0$.  Thus both structural
conclusions are also sufficient.  The second converse does not require the
additional condition $\operatorname{Re}\eta\in\Z(A^{**})$.
\end{abstract}

\noindent\fbox{\parbox{0.94\textwidth}{
\textbf{Status: candidate full result, likely valid.}
The final open question of arXiv:1907.03594 is answered affirmatively in the
same bidual-valued generality as its two structural results.  Novelty is
plausible but not certified.}}

\section{Exact source question}

Let $A$ be a $C^*$-algebra, canonically embedded in its bidual $A^{**}$.
The source proves the following one-way implications \cite{source}.

\begin{itemize}
\item Result 3.2: if a continuous linear map $\delta:A\to A^{**}$ satisfies
\[
 ab=ba=0\Longrightarrow
 a\delta(b)+\delta(a)b=b\delta(a)+\delta(b)a=0,
\]
then
\[
 \delta(a)=d(a)+\eta a,
 \qquad d:A\to A^{**}\text{ a continuous derivation},\quad
 \eta\in\Z(A^{**}).                                      \tag{R3.2}
\]
\item Result 3.3: for unital $A$, if $\delta:A\to A^{**}$ satisfies
\[
 ab^*=b^*a=0\Longrightarrow
 a\delta(b)^*+\delta(a)b^*
 =\delta(b)^*a+b^*\delta(a)=0,
\]
then there are continuous star derivations $d_1,d_2:A\to A^{**}$ and
$\eta\in A^{**}$, with $\operatorname{Re}\eta\in\Z(A^{**})$, such that
\[
 \delta(a)=d_1(a)+\eta a=d_2(a)+a\eta.                   \tag{R3.3}
\]
\end{itemize}

After proving converses under extra hypotheses in Proposition 3.4, the
source asks whether the converses of Results 3.2 and 3.3 hold in general:
\begin{center}
\includegraphics[width=0.95\textwidth]{figures/open_question_crop.png}
\end{center}

Here a star derivation is a derivation satisfying
$d(x^*)=d(x)^*$.  Products below are the canonical left and right actions
of $A$ on $A^{**}$.

\section{Idea of the proof}

\subsection*{Proof Intuition}

The requested implications are formal once one chooses the correct
decomposition.  In (R3.2), the derivation part combines by Leibniz into
$d(ab)$ or $d(ba)$, while centrality moves $\eta$ across an element of $A$;
all remaining terms contain the assumed zero product.

In (R3.3), there are two decompositions for a reason.  The expression
$a\delta(b)^*+\delta(a)b^*$ is compatible with the left-multiplier form
$d_1+\eta(\cdot)$ and collapses to $d_1(ab^*)$.  The other expression is
compatible with the right-multiplier form $d_2+(\cdot)\eta$ and collapses
to $d_2(b^*a)$.  No commutation of $\eta$ is needed.

\section{Full resolution}

\begin{theorem}\label{main}
The converses of Results 3.2 and 3.3 of \cite{source} are true.
More precisely:
\begin{enumerate}
\item Let $A$ be any $C^*$-algebra.  Every map of the form \emph{(R3.2)}
satisfies both two-sided zero-product identities of Result 3.2.
\item Let $A$ be any $C^*$-algebra for which data in \emph{(R3.3)} are
given.  Then $\delta$ satisfies both star-orthogonality identities of
Result 3.3.  This implication needs neither unitality nor
$\operatorname{Re}\eta\in\Z(A^{**})$.
\end{enumerate}
Consequently the hypotheses and structural conclusions in each source
result are equivalent (within the hypotheses under which the source proves
the forward implication).
\end{theorem}

\begin{proof}
For the first assertion, suppose $\delta(x)=d(x)+\eta x$, where $d$ is a
derivation and $\eta\in\Z(A^{**})$.  If $ab=ba=0$, associativity and
centrality give
\begin{align*}
 a\delta(b)+\delta(a)b
 &=ad(b)+d(a)b+a\eta b+\eta ab\\
 &=d(ab)+2\eta ab=0,
\end{align*}
and, with $a$ and $b$ interchanged,
\[
 b\delta(a)+\delta(b)a=d(ba)+2\eta ba=0.
\]
This proves the converse of Result 3.2.

For the second assertion, suppose
\[
 \delta(x)=d_1(x)+\eta x=d_2(x)+x\eta,
\]
where $d_1,d_2$ are star derivations.  Let $ab^*=b^*a=0$.  Use the first
representation of $\delta$ in the first required identity.  Since
$d_1(b)^*=d_1(b^*)$, we obtain
\begin{align*}
 a\delta(b)^*+\delta(a)b^*
 &=a\bigl(d_1(b)+\eta b\bigr)^*
   +\bigl(d_1(a)+\eta a\bigr)b^*\\
 &=ad_1(b^*)+d_1(a)b^*+ab^*\eta^*+\eta ab^*\\
 &=d_1(ab^*)=0.
\end{align*}
Use instead the second representation in the other required identity:
\begin{align*}
 \delta(b)^*a+b^*\delta(a)
 &=\bigl(d_2(b)+b\eta\bigr)^*a
   +b^*\bigl(d_2(a)+a\eta\bigr)\\
 &=d_2(b^*)a+b^*d_2(a)+\eta^*b^*a+b^*a\eta\\
 &=d_2(b^*a)=0.
\end{align*}
Thus Result 3.3 also has its converse.  The calculation used no property of
$\eta$ beyond belonging to the ambient bimodule algebra, proving the stated
strengthening.
\end{proof}

\section{Boundary checks and scope}

The proof is valid for nonunital $A$: it never inserts a unit or an
approximate identity.  It is also genuinely bidual-valued; every product is
taken in $A^{**}$, where the canonical copy of $A$ acts associatively.  The
centrality in (R3.2) is essential to the displayed short proof, whereas the
real-part centrality in (R3.3) is not needed for sufficiency.

There is no circular use of Proposition 3.4.  Only the definitions of a
derivation and a star derivation, the structural formulas in the conclusions,
and the stipulated zero products are used.  Thus the special assumptions in
Proposition 3.4---von Neumann or unital simple algebra and $A$-valued range---
play no role in the converses asked for.

\section{Novelty audit}

Exact-title and exact-question searches, the local run indexes, and an
OpenAlex audit of the 21 works indexed as citing the source through
17 August 2026 found no paper stating this resolution.  The citing works
concern related zero-product, Lie-centralizer, and anti-derivable-map
problems.  This supports, but cannot certify, novelty.  Because the proof is
elementary and the source wording may simply reflect an oversight, priority
is not claimed.

\begin{thebibliography}{9}
\bibitem{source}
B.~Fadaee and H.~Ghahramani,
\emph{Linear maps on $C^*$-algebras behaving like (anti-)derivations at
orthogonal elements}, Bull. Malays. Math. Sci. Soc. 43 (2020), 2851--2859;
arXiv:1907.03594.
\end{thebibliography}

\end{document}
